\documentclass[11pt]{article}
\usepackage{amsmath,amssymb,mathtools}
\usepackage[a4paper,margin=1in]{geometry}

\begin{document}

\section*{Secant Equation (0.5 pt)}

The secant equation is:
\[
B_{k+1}s_k = y_k,
\]
where \(s_k = x_{k+1}-x_k\), \(y_k=\nabla f_{k+1}-\nabla f_k\), and \(B_{k+1}\) is an approximation of the Hessian \(\nabla^2 f_{k+1}\).
\\
Prove that the secant equation is consistent with the minimization of the quadratic model problem \(m(v)\).

\paragraph{Solution.}
Consider the standard quadratic model around \(x_k\)
\[
m(v) \;=\; f(x_k) + \nabla f(x_k)^{\!\top} v + \tfrac12 v^{\!\top} B_{k+1} v ,
\qquad v\in\mathbb{R}^n.
\]
Its gradient is
\[
\nabla m(v)=\nabla f(x_k)+B_{k+1}v.
\]
To make the model consistent with the true function at the next iterate \(x_{k+1}=x_k+s_k\), we require that the model gradient interpolates the true gradient at both \(v=0\) and \(v=s_k\):
\[
\nabla m(0)=\nabla f(x_k), \qquad
\nabla m(s_k)=\nabla f(x_{k+1}).
\]
The second condition gives
\[
\nabla f(x_k) + B_{k+1} s_k \;=\; \nabla f(x_{k+1})
\quad\Longleftrightarrow\quad
B_{k+1}s_k \;=\; \nabla f(x_{k+1})-\nabla f(x_k) \;=\; y_k,
\]
which is precisely the secant equation.

\section*{Curvature Condition (0.5 pt)}

Show that the curvature condition \(s_k^{\top}y_k>0\) is required for the Quasi-Newton method.
Prove that the curvature condition is satisfied when the line search algorithm for the Wolfe conditions is used.

\paragraph{Solution.}
\textbf{(i) Derivation for a general quadratic model.}

To make the quadratic model seen before consistent with the observed change in gradient along the realized step \(s_k=x_{k+1}-x_k\), the secant equation
\(B_{k+1}s_k=y_k\) must hold.
We assume that a positive definite \(B_{k+1}\) exists, then
\[
s_k^{\top}y_k = s_k^{\top} B_{k+1} s_k > 0,
\]
since \(B_{k+1}\succ 0\) implies \(z^{\top}B_{k+1}z>0\) for all nonzero \(z\).
Therefore, \emph{if} a convex quadratic model consistent with the step exists, the curvature condition \(s_k^{\top}y_k>0\) is necessary. 

\smallskip
\textbf{(ii) Wolfe conditions.}
Let \(\Delta x\) be a descent direction at \(x\) and define
\[
\phi(t)=f(x+t\,\Delta x), \qquad \phi'(t)=\nabla f(x+t\,\Delta x)^{\!\top}\Delta x .
\]
Suppose \(t>0\) is chosen by a line search that enforces the Wolfe conditions with
\(\alpha\in(0,1)\) and \(\beta\in(\alpha,1)\):
\begin{align*}
&\text{(sufficient decrease)} & f(x+t\Delta x) &\le f(x)+ t\,\alpha\,\nabla f(x)^{\!\top}\Delta x &&\iff&& \phi(t)\le \phi(0)+ t\,\alpha\,\phi'(0),\\
&\text{(curvature)} & \nabla f(x+t\Delta x)^{\!\top}\Delta x &\ge \beta\,\nabla f(x)^{\!\top}\Delta x &&\iff&& \phi'(t)\ge \beta\,\phi'(0).
\end{align*}
Set \(s_k=t\,\Delta x\) and \(y_k=\nabla f(x+t\Delta x)-\nabla f(x)\).
Then
\[
s_k^{\top}y_k
= t\Big(\Delta x^{\!\top}\nabla f(x+t\Delta x)-\Delta x^{\!\top}\nabla f(x)\Big)
= t\big(\phi'(t)-\phi'(0)\big).
\]
By the curvature Wolfe condition,
\(\phi'(t)\ge \beta\,\phi'(0)\).
Since \(\Delta x\) is a descent direction, \(\phi'(0)=\nabla f(x)^{\!\top}\Delta x<0\).
Therefore,
\[
s_k^{\top}y_k \;\ge\; t\big(\beta\,\phi'(0)-\phi'(0)\big)
= t(\beta-1)\,\phi'(0) \;>\; 0,
\]
because \(\beta-1<0\), \(\phi'(0)<0\), and \(t>0\).
Hence any step length chosen under these Wolfe conditions satisfies the curvature condition \(s_k^{\top}y_k>0\).

\end{document}
